% Part: many-valued-logic
% Chapter: three-valued-logics
% Section: goedel

\documentclass[../../../include/open-logic-section]{subfiles}

\begin{document}

\olfileid{mvl}{thr}{god}

\olsection{منطقات غودل}

وضع كورت غودل سلسلة من المنطقات ذات~$n$ من القيم؛ يشتمل كل واحد
منها على جميع~!!{formula}s الصحيحة حدسيًا، ويشتمل عليه المنطق
الكلاسيكي. وهذا أول أفراد السلسلة غير التافهة:

\begin{defn}\ollabel{defn:goedel}
مصفوفة \emph{منطق غودل الثلاثي القيم}~$\LogGod$ هي:
\begin{enumerate}
  \item لغة القضايا القياسية~$\Lang L_0$ ذات الروابط
  $\lfalse$ و$\lnot$ و$\land$ و$\lor$ و$\lif$.
  \item مجموعة قيم الصدق~$V = \{\True, \Undef, \False\}$.
  \item القيمة المميزة وحدها هي~$\True$؛ أي~$V^+ = \{\True\}$.
  \item في الثابت~$\lfalse$ يكون~$\tf{\lfalse} = \False$؛ وأما
  دوال صدق سائر الروابط فتعطيها الجداول:
  \begin{center}
    \begin{tabular}{c|c}
      $\tf{\lnot}[\LogGod]$ & \\
      \hline
      $\True$ & $\False$ \\
      $\Undef$ & $\False$ \\
      $\False$ & $\True$
    \end{tabular}
    \quad
    \begin{tabular}{c|ccc}
      $\tf{\land}[\LogGod]$ & $\True$ & $\Undef$ & $\False$ \\
      \hline
      $\True$ & $\True$ & $\Undef$ & $\False$ \\
      $\Undef$ & $\Undef$ & $\Undef$ & $\False$\\
      $\False$ & $\False$ & $\False$ & $\False$
    \end{tabular}
    \\[2ex]
    \begin{tabular}{c|ccc}
      $\tf{\lor}[\LogGod]$ & $\True$ & $\Undef$ & $\False$ \\
      \hline
      $\True$ & $\True$ & $\True$ & $\True$ \\
      $\Undef$ & $\True$ & $\Undef$ & $\Undef$ \\
      $\False$ & $\True$ & $\Undef$ & $\False$
    \end{tabular}
    \quad
    \begin{tabular}{c|ccc}
      $\tf{\lif}[\LogGod]$ & $\True$ & $\Undef$ & $\False$ \\
      \hline
      $\True$ & $\True$ & $\Undef$ & $\False$ \\
      $\Undef$ & $\True$ & $\True$ & $\False$  \\
      $\False$ & $\True$ & $\True$ & $\True$
    \end{tabular}
  \end{center}
\end{enumerate}
\end{defn}

وجدولا~$\land$ و$\lor$ هنا هما الجدولان نفسيهما في منطق~\L
ukasiewicz وفي منطق كليني القوي؛ وإنما يختلف~$\lnot$ و$\lif$ في
هذه المنطقات. ففي منطق غودل~$\tf{\lnot}(\Undef) = \False$.
وكذلك، بخلاف منطق~\L ukasiewicz ومنطق كليني، يكون
$\tf{\lif}(\Undef, \False) = \False$؛ وأما
$\tf{\lif}(\Undef, \Undef) = \True$، فيخالف منطق كليني ويوافق
منطق~\L ukasiewicz.

وارتباط هذا المنطق بالمنطق الحدسي ظاهر في~$\LogGod[3]$: فكل حقيقة
حدسية تحصيل في~$\LogGod[3]$، وكثير من التحصيلات الكلاسيكية التي لا تصح حدسيًا لا
تكون تحصيلات في~$\LogGod[3]$ أيضًا. فمن أمثلتها هذه الصيغ:
\begin{align*}
  & p \lor \lnot p && (p \lif q) \lif (\lnot p \lor q) \\
  & \lnot\lnot p \lif p && \lnot(\lnot p \land \lnot q) \lif (p \lor q) \\
  & ((p \lif q) \lif p) \lif p && \lnot(p \lif q) \lif (p \land \lnot q)
\end{align*}
ومع ذلك توجد تحصيلات لـ~$\LogGod[3]$ لا تصح حدسيًا، ومنها
$\lnot\lnot p \lor \lnot p$ و~$(p \lif q) \lor (q \lif p)$.

\begin{prob}
  ضع جداول صدق تثبت أن الصيغ الآتية تحصيلات في~$\LogGod[3]$:
  \begin{align*}
    & \lnot\lnot p \lor \lnot p\\
    & (p \lif q) \lor (q \lif p) \\
    & \lnot(p \land q) \lif (\lnot p \lor \lnot q) \\
    & (p \lif q) \lor (q \lif r) \lor (r \lif s)
  \end{align*}
\end{prob}

\begin{prob}
  ضع جداول صدق تثبت أن الصيغ الآتية ليست تحصيلات في~$\LogGod[3]$:
  % تصحيح أسلوبي (0396-N1): استُعمل فعل الأمر بدل المصدر.
  \begin{align*}
    & (p \lif q) \lif (\lnot p \lor q) \\
    & \lnot(\lnot p \land \lnot q) \lif (p \lor q) \\
    & ((p \lif q) \lif p) \lif p \\
    & \lnot(p \lif q) \lif (p \land \lnot q)
  \end{align*}
\end{prob}

\begin{prob}
  أي العلاقات الآتية تصح في منطق غودل؟ ضع جدول صدق لكل واحدة.
  \begin{enumerate}
    \item $p, p \lif q \Entails q$
    \item $p \lor q, \lnot p \Entails q$
    \item $p \land q \Entails p$
    \item $p \Entails p \land p$
    \item $p \Entails p \lor q$
  \end{enumerate}
\end{prob}

\end{document}
